\section{Méthodologie}

\subsection{Vue d'ensemble}

Le cadre proposé est une chaîne de traitement complète pour la détection
automatique des fuites à partir de séries temporelles multivariées de capteurs.
Il se compose de quatre étapes principales : (i) prétraitement des données et
ingénierie de caractéristiques, (ii) construction de séquences par approche
fenêtre glissante, (iii) classification via une architecture hybride
CNN-BiLSTM-Attention, et (iv) calibration du seuil pour la prise de décision
opérationnelle.

La même stratégie de prétraitement et de découpage des données est appliquée à
tous les modèles considérés dans cette étude afin de garantir une comparaison
équitable dans l'évaluation de benchmark.

\subsection{Formulation par fenêtre glissante}

Soit une série temporelle multivariée $\mathbf{X} = \{x_1, x_2, \dots, x_T\}$,
où $x_t \in \mathbb{R}^d$ représente le vecteur de caractéristiques à l'instant
$t$. Les données sont transformées en échantillons d'apprentissage supervisé à
l'aide d'une fenêtre glissante de taille $W$ et de pas $S$.

Chaque séquence d'entrée est définie par :

\begin{equation}
\mathbf{X}_t = \{x_t, x_{t+1}, \dots, x_{t+W-1}\}
\end{equation}

avec l'étiquette associée :

\begin{equation}
y_t \in \{0,1\}
\end{equation}

où $y_t = 1$ indique la présence d'au moins un événement de fuite dans la fenêtre.

Dans ce travail, une taille de fenêtre $W = 20$ et un pas $S = 5$ sont utilisés,
ce qui permet au modèle de capturer des motifs temporels de court terme tout en
préservant un recouvrement suffisant entre séquences consécutives.

\subsection{Ingénierie de caractéristiques}

Pour enrichir la représentation des signaux bruts des capteurs, des
caractéristiques guidées par le domaine sont construites, notamment :

\begin{itemize}
    \item Le gradient de pression ($\Delta P_t$) pour capturer les variations brusques de pression,
    \item La normalisation du débit (z-score) pour mettre en évidence les anomalies relatives,
    \item Le ratio pression/débit pour modéliser les relations hydrauliques.
\end{itemize}

Ces caractéristiques sont concaténées aux variables d'origine pour former des
vecteurs de caractéristiques enrichis à chaque pas de temps.

\subsection{Architecture du modèle}

Le modèle proposé intègre des composantes convolutionnelles, récurrentes et
basées sur l'attention afin de capturer à la fois les dépendances temporelles
locales et globales.

\subsubsection{Couches convolutionnelles}

Les séquences d'entrée sont d'abord traitées par des couches convolutionnelles
unidimensionnelles afin d'extraire les motifs temporels locaux :

\begin{equation}
h^{(c)} = \text{ReLU}(\text{Conv1D}(\mathbf{X}))
\end{equation}

Ces couches agissent comme des extracteurs de caractéristiques, en capturant des
variations de court terme telles que les chutes de pression et les pics de débit.

\subsubsection{LSTM bidirectionnel}

Les caractéristiques extraites sont ensuite passées à un LSTM bidirectionnel
(BiLSTM), qui capture les dépendances de long terme dans les directions avant
et arrière :

\begin{equation}
h^{(l)} = \text{BiLSTM}(h^{(c)})
\end{equation}

Cela permet au modèle d'intégrer l'information contextuelle provenant à la fois
des pas de temps passés et futurs dans chaque séquence.

\subsubsection{Mécanisme d'auto-attention}

Pour mettre l'accent sur les pas de temps les plus informatifs, un mécanisme
d'auto-attention est appliqué aux sorties du BiLSTM :

\begin{equation}
\alpha_t = \frac{\exp(e_t)}{\sum_{k} \exp(e_k)}, \quad
e_t = \mathbf{v}^\top \tanh(\mathbf{W} h^{(l)}_t)
\end{equation}

Le vecteur de contexte résultant est calculé comme suit :

\begin{equation}
c = \sum_t \alpha_t h^{(l)}_t
\end{equation}

Ce mécanisme permet au modèle de se concentrer sur des signatures de fuite
temporellement localisées tout en préservant l'information globale de la séquence.

\subsubsection{Couche de sortie}

Le vecteur de contexte passe ensuite par une couche entièrement connectée avec
activation sigmoïde pour produire la prédiction finale :

\begin{equation}
\hat{y} = \sigma(\mathbf{W}_o c + b)
\end{equation}

où $\hat{y} \in [0,1]$ représente la probabilité de fuite.

\subsection{Stratégie d'entraînement}

Le modèle est entraîné avec la fonction de perte entropie croisée binaire :

\begin{equation}
\mathcal{L} = - \frac{1}{N} \sum_{i=1}^{N} \left[
y_i \log(\hat{y}_i) + (1 - y_i)\log(1 - \hat{y}_i)
\right]
\end{equation}

Le jeu de données est découpé en ensembles d'entraînement, de calibration et de
test à l'aide d'une séparation temporelle pour éviter les fuites de données.
L'arrêt anticipé est appliqué selon l'AUC-PR de validation, et une stratégie
d'ajustement du taux d'apprentissage est utilisée pour améliorer la convergence.

Tous les modèles de référence sont entraînés avec les mêmes découpages de données
et la même chaîne de prétraitement afin de garantir une comparaison de benchmark
équitable.

\subsection{Calibration du seuil}

Au lieu d'utiliser un seuil fixe (par exemple 0.5), le seuil de décision est
optimisé sur un ensemble de calibration pour équilibrer précision et rappel :

\begin{equation}
\hat{y}_{final} =
\begin{cases}
1 & \text{if } \hat{y} \geq \tau \\
0 & \text{otherwise}
\end{cases}
\end{equation}

Cela permet au modèle d'opérer selon différents compromis en fonction des
exigences applicatives, comme la priorité à la sensibilité de détection ou la
minimisation des fausses alertes.

\subsection{Détails d'implémentation}

Le modèle est implémenté à l'aide d'un framework de deep learning avec une
architecture compacte comprenant environ 23{,}000 paramètres entraînables.
L'optimisation est effectuée avec l'optimiseur Adam, et l'entraînement est
réalisé par descente de gradient en mini-lots sur plusieurs époques.

Toutes les expériences sont menées dans des conditions cohérentes afin d'assurer
la reproductibilité et une comparaison équitable entre les différents modèles.